%% 11장 초월함수를 만드는 법
\chapter{Implementing Elementary Functions}
\label{ch:elementary-functions}

\abstract{컴퓨터의 산술 장치가 하는 일은 덧셈과 곱셈과 나눗셈뿐인데, 우리는 \code{exp(x)}나 \code{cos(x)}를 부르면 마지막 비트 가까이 정확한 값이 나온다고 당연하게 여긴다. 이 장은 그 값이 어떻게 만들어지는지를 다룬다. 기본 전략은 셋이다. 인자를 좁은 구간으로 옮기는 인자 축소, 그 좁은 구간에서 함수를 다항식으로 근사하는 것, 그리고 그 다항식을 빠르고 정확하게 계산하는 것이다. 인자 축소에서는 상수를 쪼개 곱셈을 정확하게 만드는 Cody--Waite의 방법과 표를 쓰는 Tang의 방법을, 근사에서는 Taylor 다항식이 최선이 아닌 이유와 Chebyshev의 교대 정리, 최선 근사를 찾는 Remez의 알고리즘을, 계산에서는 Horner의 방법과 병렬성을 살린 Estrin의 방법을 본다. 마지막으로 결과를 올바르게 반올림하는 것이 왜 어려운지를 다룬다. 예로 삼는 것은 stardust가 직접 만든 \code{exp}, \code{expm1}, \code{log1p}, \code{cos} 커널이다.}

\section{A Three-Step Strategy}
\label{sec:three-steps}

$e^x$를 계산한다고 하자. Taylor 급수 $\sum x^k/k!$은 모든 $x$에서 수렴하지만 $x=700$에서는 항이 $700^{700}/700!$ 근처까지 커졌다가 줄어들므로 수천 개의 항이 필요하고, 음수 $x$에서는 큰 항들이 서로 지워져 결과가 엉망이 된다(9장의 소거). 급수는 $x$가 $0$ 근처일 때만 쓸모가 있다. 그러니 먼저 $x$를 $0$ 근처의 작은 수로 옮겨야 한다. 지수함수에는 그렇게 하기 좋은 성질 $e^{a+b}=e^ae^b$가 있다. $x=k\ln2+r$이고 $|r|\le\ln2/2$이면 $e^x=2^ke^r$이며, $2^k$를 곱하는 것은 지수에 $k$를 더하는 정확한 연산이다(8장). 남은 $e^r$은 좁은 구간에서 다항식 하나로 근사할 수 있다.

이것이 거의 모든 초월함수 구현의 틀이다. 첫째로 함수의 대수적 성질을 써서 인자를 좁은 \term{핵심 구간}{core interval}으로 옮기고(\term{인자 축소}{argument reduction}), 둘째로 그 구간에서 함수를 다항식이나 유리함수로 근사하며, 셋째로 축소의 역을 적용해 결과를 재구성한다. 각 단계의 오차가 최종 결과에 어떻게 전해지는지를 따져야 하는데 마지막 비트까지 맞추려면 셋 모두의 정확도가 반올림 한 번에 가까워야 하므로 10장의 오차 없는 변환이 곳곳에서 쓰인다.

\section{Argument Reduction}
\label{sec:reduction}

\subsection{Exact Reduction: Cody--Waite}

$x=k\ln2+r$에서 $k=\operatorname{round}(x/\ln2)$로 정하고 $r=x-k\ln2$를 계산한다. 그런데 $\ln2$는 binary64 수가 아니다. $\ln2$를 반올림한 값 $C$를 쓰면 $kC$와 $k\ln2$의 차이가 $|k|\cdot u\ln2$ 정도이고, $r$이 작으므로 이 차이가 $r$에 대한 상대오차로는 크게 부풀려진다. 9장의 언어로는 뺄셈 $x-k\ln2$의 조건수가 $|x|/|r|$로 크다.

Cody와 Waite의 해법은 상수를 두 조각으로 쪼개는 것이다\cite{cody1980software}. $\ln2=C_1+C_2+\cdots$로 쓰되 $C_1$의 유효숫자 아래쪽 몇 비트를 $0$으로 두면, $k$가 작은 정수일 때 곱 $kC_1$이 $53$비트 안에 들어가 정확하다. 그러면 $x-kC_1$은 $x$와 $kC_1$이 가까우므로 Sterbenz의 보조정리에서 정확하고, 남은 $-kC_2$를 빼는 단계에서만 반올림이 일어난다. FMA가 있으면 더 간단하다. $\fl(x-kC_1)$을 FMA 한 번으로 계산하면 $x-kC_1$이 정확히 계산된 뒤 한 번 반올림되는데, 결과가 $\mathbb F$에 있으면 반올림이 아무것도 바꾸지 않는다.

$k$를 구하는 방법에도 비트 수준의 기교가 있다. $1.5\cdot2^{52}$를 더하면 된다. $|y|<2^{51}$인 $y$에 이 수를 더하면 결과가 이진 구간 $[2^{52},2^{53})$에 들어가는데 그 구간의 간격은 $1$이므로 반올림이 곧 가장 가까운 정수로 반올림하는 것이고 다시 $1.5\cdot2^{52}$를 빼면 Sterbenz의 보조정리에서 정확하게 $\operatorname{round}(y)$가 남는다. 게다가 더한 결과의 비트 패턴에서 아래쪽 비트들이 곧 그 정수의 2의 보수 표현이어서, 정수 변환 명령 없이 정수 $k$를 꺼낼 수 있다.

\subsection{Table-Based Reduction}

$|r|\le\ln2/2\approx0.35$는 아직 넓다. 이 구간에서 $e^r$을 binary64 정밀도로 근사하려면 다항식의 차수가 $11$쯤 필요하다. Tang은 축소를 한 단계 더 하되 표를 쓰는 방법을 제안했다\cite{tang1989exp}. $N=2^m$을 정하고 $x=(k+j/N)\ln2+r$, $0\le j<N$, $|r|\le\ln2/(2N)$로 쓰면
\[
e^x=2^k\cdot2^{j/N}\cdot e^r
\]
이고, $2^{j/N}$의 $N$개 값을 미리 높은 정밀도로 계산해 표로 둔다. 남은 $e^r$은 훨씬 좁은 구간에서 근사하므로 낮은 차수로 충분하다. 표 하나에는 $2^{j/N}$을 binary64로 반올림한 값과 그 반올림 오차인 상대 꼬리를 함께 담아서 표의 값이 사실상 double-double만큼 정밀해진다.

로그는 곱의 성질 $\log(ab)=\log a+\log b$로 축소한다. $x=2^kz$로 지수를 떼어 내면 $\log x=k\log2+\log z$이고, $z$가 속한 작은 구간마다 대표값 $c$와 그 역수 $1/c$의 근삿값, $\log c$의 높은 정밀도 값을 표로 두면 $\log z=\log c+\log(z/c)$에서 $z/c$가 $1$에 아주 가까워진다. 이때 $z/c-1$을 나눗셈 없이 FMA 한 번 $r=z\cdot(1/c)-1$로 계산하는 것이 요점이다. 표의 $1/c$가 정확한 역수가 아니어도 괜찮은데, $\log c$를 그 $1/c$에 맞추어, 곧 표의 $1/c$의 정확한 역수를 $c$로 삼아 계산해 두면 되기 때문이다.

\begin{instardust}[\code{csrc/array/kernels/vmath.hpp}]
배열의 \code{exp} 커널은 $N=128$인 Tang의 축소이다. $z=x\cdot(128/\ln2)+1.5\cdot2^{52}$로 정수 $128k+j$를 얻고, $\ln2/128$을 두 조각으로 쪼갠 상수로 FMA 두 번을 써서 $r$을 만든다. 표의 $j$번째 칸은 $2^{j/128}$의 binary64 값과 상대 꼬리를 Arb로 계산해 반올림한 것이며, $z$의 비트 패턴을 $45$비트 왼쪽으로 밀어 표 값의 비트에 정수로 더하면 지수 필드에 $k$가 더해져 $2^k$를 곱한 셈이 된다($45=52-7$이고, $2^7=128$이다). \code{log} 핵심은 대략 $z\in[0.686,1.371)$을 비트 패턴으로 $128$개 구간에 나누고, $1$을 담은 구간은 $1$이 가운데에 오도록 해서 $c=1$로 둔다. 그러면 $x$가 $1$ 근처일 때 축소가 아무 오차도 만들지 않는다.
\end{instardust}

\subsection{Reduction for Trigonometric Functions}

$\cos x$는 $x=n\pi/2+r$, $|r|\le\pi/4$로 축소하고 $n\bmod4$에 따라 $\pm\cos r$이나 $\pm\sin r$을 고른다. 원리는 지수함수와 같지만 두 가지가 더 어렵다. 하나는 $\pi/2$를 더 많은 조각으로 쪼개야 한다는 것이다. 삼각함수에서는 $r$ 자체가 결과의 크기를 정하므로 $r$의 상대 정확도가 그대로 결과의 정확도가 되고, $x$가 $\pi/2$의 정수배에 아주 가까우면 $r$이 극도로 작아질 수 있기 때문이다. 다른 하나는 $x$가 아주 클 수 있다는 것이다. $x=10^{22}$이면 $n$이 $10^{22}$ 근처이고, $\pi/2$를 약 $200$비트까지 알아야 $r$을 정확히 얻고, binary64의 모든 인자를 덮으려면 천 비트 남짓이 필요하다. 이런 인자를 위해 Payne과 Hanek은 $2/\pi$의 이진 전개를 표로 두고 $x$의 유효숫자와 곱하되 결과에 영향을 주는 비트만 골라 곱하는 방법을 만들었다\cite{payne1983}.

\begin{instardust}[\code{csrc/array/kernels/vmath.hpp}]
배열의 \code{cos} 커널은 $|x|\le2^{20}$에서 Cody--Waite 축소를 쓴다. $\pi/2$를 네 조각 $P_1+P_2+P_3+P_4$로 쪼개는데, $P_1$과 $P_2$는 $33$비트로 잘라 두었으므로 $n<2^{20}$이면 $nP_1$과 $nP_2$가 $53$비트 안에 들어가 정확하다. 둘째 뺄셈의 반올림 오차도 FMA로 정확히 구해 두고 축소 결과를 double-double $r_{\mathrm{hi}}+r_{\mathrm{lo}}$로 넘긴다. $2^{20}$보다 큰 인자는 Payne--Hanek 축소를 갖춘 C 라이브러리의 \code{cos}에 맡긴다.
\end{instardust}

\section{Polynomial Approximation}
\label{sec:poly-approx}

\subsection{Taylor Polynomials Are Not Optimal}

핵심 구간 $[-a,a]$에서 $e^r$을 $n$차 다항식으로 근사한다고 하자. 가장 먼저 떠오르는 것은 Taylor 다항식 $\sum_{k\le n}r^k/k!$이고, 그 오차는 대략 $a^{n+1}/(n+1)!$이다. stardust의 \code{exp} 커널에서 $a=\ln2/256\approx2.71\times10^{-3}$이므로 $n=5$이면 오차가 약 $5.47\times10^{-19}\approx2^{-60.7}$이다. $e^r\approx1$이므로 이것이 곧 상대오차이고, 결과의 반올림 오차 $2^{-53}$에 비하면 충분히 작다.

그런데 Taylor 다항식은 $0$ 한 점에서 함수에 가장 잘 맞도록 만든 것이어서, 오차가 구간의 끝에 몰려 있고 가운데에서는 쓸데없이 작다. 구간 전체에서 가장 큰 오차를 최소로 하는 다항식, 곧 \term{최선 근사}{best approximation} 또는 \term{미니맥스 근사}{minimax approximation}를 쓰면 같은 차수로 오차를 크게 줄일 수 있다. 위의 예에서 $5$차 미니맥스 다항식의 최대 오차는 약 $1.71\times10^{-20}$으로 Taylor의 $1/32$이다. 이 비 $32=2^5$은 우연이 아니며, 다음 절의 Chebyshev 다항식이 그 이유를 설명한다. 이 수치들은 모두 십진 $60$자리 산술로 직접 계산했다.

차수를 하나 낮추면 어떨까. $4$차 미니맥스 다항식의 오차는 $7.6\times10^{-17}\approx2^{-53.6}$으로 단위 반올림과 비슷해 모자란다. 그래서 이 구간에서는 $5$차가 필요하고, $5$차라면 Taylor 다항식으로도 충분하다. 계수가 $1/k!$이라는 단순한 꼴이어서 반올림 오차가 적고 설명하기 쉽다는 장점까지 있으므로, 이 경우에는 Taylor 다항식이 합리적인 선택이다. 근사의 차수와 종류는 결국 표의 크기, 다항식의 차수, 계수의 반올림 사이의 균형으로 정해진다.

\subsection{Chebyshev's Alternation Theorem}

최선 근사를 특징짓는 것은 오차의 모양이다.

\begin{theorem}[Chebyshev의 교대 정리]
\label{thm:equioscillation}
$f$가 $[a,b]$에서 연속이고 $\mathcal P_n$이 차수 $n$ 이하의 실계수 다항식 전체일 때, $p\in\mathcal P_n$이 $\max_{[a,b]}|f-p|$를 최소로 하는 필요충분조건은 $a\le x_0<x_1<\cdots<x_{n+1}\le b$인 점 $n+2$개가 있어서
\[
f(x_i)-p(x_i)=\sigma(-1)^iE,\qquad E=\max_{[a,b]}|f-p|,\quad\sigma\in\{1,-1\}
\]
가 성립하는 것이다. 이런 $p$는 유일하다.
\end{theorem}

\begin{proof}
충분성만 증명한다. 오차가 그런 $n+2$개의 점에서 교대로 $\pm E$를 갖는데, 어떤 $q\in\mathcal P_n$이 $\max|f-q|<E$를 만족한다고 하자. 각 $x_i$에서
\[
q(x_i)-p(x_i)=\bigl(f(x_i)-p(x_i)\bigr)-\bigl(f(x_i)-q(x_i)\bigr)
\]
의 첫 괄호는 크기가 $E$이고 둘째 괄호는 크기가 $E$보다 작으므로, $q-p$의 부호는 $f-p$의 부호와 같다. 따라서 $q-p$는 $n+2$개의 점에서 부호를 번갈아 바꾸고, 중간값 정리에서 근을 적어도 $n+1$개 갖는다. $q-p\in\mathcal P_n$이므로 $q=p$인데, 그러면 $\max|f-q|=E$여서 가정과 모순이다. 필요성과 유일성의 증명은 \cite{cheney1998}에 있다.
\end{proof}

이 정리에서 $[-1,1]$의 함수 $x^{n+1}$을 차수 $n$ 이하의 다항식으로 근사하는 문제를 풀 수 있다. Chebyshev 다항식 $T_{n+1}(x)=\cos((n+1)\arccos x)$는 최고차항 계수가 $2^n$이고 $[-1,1]$에서 $n+2$개의 점 $\cos(k\pi/(n+1))$에서 교대로 $\pm1$을 가지므로, $x^{n+1}-2^{-n}T_{n+1}(x)$가 $x^{n+1}$의 최선 근사이고 오차의 최댓값은 $2^{-n}$이다. 매끄러운 함수의 Taylor 나머지는 대략 다음 항 $c\,x^{n+1}$이므로, 이것을 Chebyshev 다항식으로 바꾸어 빼면 오차가 대략 $2^{-n}$배로 준다. $5$차 근사에서 그 비는 $2^{-5}=1/32$로 위의 측정과 정확히 맞는다. 미니맥스 다항식은 이 교정을 모든 차수에서 최적으로 한 것이라고 볼 수 있다.

\subsection{The Remez Algorithm}

교대 정리는 최선 근사를 계산하는 방법도 알려 준다. 교대점 $x_0<\cdots<x_{n+1}$을 안다면, 미지수 $c_0,\dots,c_n,E$에 대한 일차연립방정식
\[
\sum_{j=0}^nc_jx_i^j+(-1)^iE=f(x_i),\qquad i=0,\dots,n+1
\]
을 풀어 최선 근사를 얻는다. 교대점을 모르므로 Remez는 다음과 같이 반복한다. 처음에는 Chebyshev 점을 교대점으로 삼아 위의 연립방정식을 풀고, 얻은 다항식의 오차 함수 $f-p$의 국소 극값들을 찾아 새로운 교대점으로 삼아 다시 푼다. 이 \term{교환 알고리즘}{exchange algorithm}은 적당한 조건에서 이차로 수렴한다\cite{trefethen2019atap}.

\begin{algorithm}[t]
\begin{algorithmic}[1]
\Require 연속함수 $f$, 구간 $[a,b]$, 차수 $n$
\Ensure 미니맥스 다항식 $p$와 오차의 크기 $|E|$
\State $x_i\gets(a+b)/2-((b-a)/2)\cos(i\pi/(n+1))$, $i=0,\dots,n+1$ \Comment{Chebyshev 점}
\Repeat
  \State $\sum_jc_jx_i^j+(-1)^iE=f(x_i)$를 풀어 $p=\sum_jc_jx^j$와 $E$를 얻는다
  \State $f-p$의 국소 극값 $n+2$개를 찾아 새 $x_i$로 삼는다 \Comment{양 끝점 포함}
\Until{극값의 크기가 모두 $|E|$와 충분히 가깝다}
\State \Return $p$, $|E|$
\end{algorithmic}
\caption{Remez의 교환 알고리즘}
\label{alg:remez}
\end{algorithm}

실제 라이브러리를 만들 때는 두 가지를 더 고려한다. 첫째로 계수는 binary64로 반올림되어야 하므로, 반올림한 계수로도 오차가 작게 유지되는 다항식을 찾아야 한다. 둘째로 절대오차보다 상대오차를 최소로 해야 하는 경우가 많은데, 이것은 가중치 $1/|f|$를 준 미니맥스 문제가 되고 Remez 알고리즘을 조금 고쳐 풀 수 있다. 차수가 같은 다항식이라도 Chebyshev 점에서 보간한 다항식은 미니맥스 다항식에 아주 가까운 오차를 내는데 위의 예에서 $5$차 Chebyshev 보간의 오차는 $1.7108\times10^{-20}$으로 미니맥스의 $1.7101\times10^{-20}$과 거의 같다. 보간 이론은 29장에서 다룬다.

\section{Polynomial Evaluation}
\label{sec:poly-eval}

\subsection{Horner's Method}

$p(x)=\sum_{k=0}^na_kx^k$를 계산하는 표준적인 방법은 괄호를 안에서부터 묶는
\[
p(x)=a_0+x\bigl(a_1+x\bigl(a_2+\cdots+x(a_{n-1}+xa_n)\cdots\bigr)\bigr)
\]
꼴인 \term{Horner의 방법}{Horner's method}이다. 곱셈과 덧셈이 $n$번씩이면 되고, FMA를 쓰면 FMA $n$번이다. 반올림 오차는 9장의 방법으로 분석하면
\[
|\hat p(x)-p(x)|\le\gamma_{2n}\sum_{k=0}^n|a_k||x|^k
\]
이다. 계수 $a_k$와 거듭제곱 $x^k$가 각각 많아야 $2n$번의 반올림을 거치기 때문이다. 우변의 합은 계수의 부호가 섞이면 $|p(x)|$보다 클 수 있지만 핵심 구간의 근사 다항식은 대개 첫 항이 나머지를 압도하므로($e^r\approx1+r+\cdots$에서 $1$이 지배한다) 상대오차가 $u$의 몇 배 수준으로 유지된다.

\subsection{Estrin's Method}

Horner의 방법에는 속도의 약점이 있다. 각 단계가 앞 단계의 결과를 기다려야 하므로, 연산 하나의 지연 시간이 $L$사이클이면 전체에 $nL$사이클이 걸린다. 현대 프로세서는 서로 독립인 연산 여러 개를 동시에 처리할 수 있으므로(13장), 의존 관계의 사슬을 짧게 만들면 빨라진다. Estrin은 이웃한 두 계수씩 일차식으로 묶고 그것들을 $x^2$, $x^4$, $\dots$로 다시 짝지어
\[
p(x)=(a_0+a_1x)+x^2(a_2+a_3x)+x^4\bigl((a_4+a_5x)+x^2(a_6+a_7x)\bigr)
\]
처럼 묶었다\cite{estrin1960}. 괄호 안의 일차식들은 서로 독립이어서 동시에 계산할 수 있고, $x^2$과 $x^4$도 따로 만들 수 있으므로 의존 사슬의 길이가 $n$에서 대략 $\log_2n$으로 줄어든다. 대신 곱셈이 몇 번 늘고 반올림 오차의 상한이 조금 커진다.

\begin{instardust}[\code{csrc/array/kernels/vmath.hpp}]
\code{exp} 커널의 다항식은 짧은 Estrin 꼴이다. $t=\mathit{tail}+r+r^2(1/2+r/6)+r^4(1/24+r/120)$을
\[
p_{23}=\operatorname{fma}(r,1/6,1/2),\quad p_{45}=\operatorname{fma}(r,1/120,1/24),\quad t=\operatorname{fma}(r^4,p_{45},\operatorname{fma}(r^2,p_{23},\mathit{tail}+r))
\]
로 계산하므로 $p_{23}$, $p_{45}$, $r^2$이 동시에 계산된다. 마지막에 $e^x=s(1+t)$를 $\operatorname{fma}(s,t,s)$로 만들어 반올림을 한 번으로 줄인다. 여기서 $s=2^{k+j/128}$은 표의 값이고 $\mathit{tail}$은 그 상대 꼬리이다. \code{cos} 커널은 FreeBSD의 \code{\_\_kernel\_cos}와 \code{\_\_kernel\_sin}을 따르는데, 계수는 fdlibm의 미니맥스 계수이고 근사 오차가 상대적으로 $2^{-58}$ 이하이다.
\end{instardust}

\section{Correct Rounding}
\label{sec:correct-rounding}

\subsection{The Table Maker's Dilemma}

8장에서 보았듯이 IEEE~754는 기본 연산에 올바른 반올림을 요구하지만 초월함수에는 요구하지 않는다. 그 이유는 다음과 같다. 근삿값 $\hat y$가 참값 $y=f(x)$의 상대오차 $\epsilon$ 안에 있다는 것을 안다고 하자. $\hat y$를 반올림한 결과가 $y$를 반올림한 결과와 같다고 확신하려면, 구간 $[\hat y(1-\epsilon),\hat y(1+\epsilon)]$ 안에 반올림의 경계, 곧 이웃한 두 binary64 수의 한가운데가 없어야 한다. 참값이 경계에 아주 가까우면 경계를 피할 만큼 $\epsilon$을 작게 하려면 정밀도를 아주 높여야 한다. 문제는 참값이 경계에 얼마나 가까울 수 있는지 미리 알기 어렵다는 것이다. 수표를 만들던 시절부터 알려진 이 문제를 \term{표 만드는 사람의 딜레마}{table maker's dilemma}라 한다.

확률적으로 생각하면 함수값의 비트들은 무작위처럼 보이므로, 반올림 경계에서 $2^{-53-m}$ 이내에 있는 참값은 대략 $2^{-m}$의 비율로 나타난다. binary64의 인자는 $2^{64}$개 정도이므로 가장 어려운 인자는 경계에서 $2^{-53-64}$ 정도까지 가까울 것으로 예상되고, 그렇다면 약 $120$비트 정밀도로 계산하면 언제나 올바르게 반올림할 수 있을 것이다. 이 어림은 실제로 가장 어려운 인자들을 모두 찾아내는 계산으로 확인되었다\cite{lefevre1998}.

\subsection{Ziv's Strategy}

그렇다고 모든 계산을 $120$비트로 할 필요는 없다. Ziv는 먼저 binary64보다 조금 높은 정밀도로 빠르게 계산하고, 그 결과의 오차 구간이 반올림 경계를 포함하지 않으면 바로 반올림해 끝내고, 포함하면 더 높은 정밀도로 다시 계산하는 전략을 제안했다\cite{ziv1991}. 대부분의 인자는 첫 단계에서 끝나므로 평균 속도는 빠르고, 드문 어려운 인자에서만 느려진다. 10장에서 본 Arb로 참값을 반올림하는 방법도 같은 생각이다. 최근의 CORE-MATH 프로젝트는 이 전략과 가장 어려운 인자의 목록을 결합해 binary64의 모든 기본 초월함수를 올바르게 반올림하면서도 빠르게 계산하려는 시도이다\cite{coremath2022}.

일반적인 수학 라이브러리는 올바른 반올림을 약속하지 않고 대신 오차가 $1$ ULP 이내라든가 하는 상한을 약속한다. 이 경우 대부분의 인자에서는 올바르게 반올림된 값이 나오지만 일부에서는 이웃한 수가 나온다. 그래서 같은 함수라도 라이브러리마다 일부 인자에서 마지막 비트가 다르고, 이것이 서로 다른 시스템에서 같은 계산이 다른 비트를 내는 흔한 원인이다.

\begin{instardust}[\code{csrc/array/kernels/vmath.hpp}]
파일의 주석은 각 커널을 Arb의 참값과 비교한 결과를 함께 적어 둔다. 함수마다 전 범위의 무작위 인자와 $0$, $1$, $\pi/2$의 배수 근처에 몰아 둔 인자 $10$만--$20$만 개에서 잰 것으로, 다섯 커널 모두 최대 오차가 $1$ ULP이고 올바르게 반올림되지 않은 인자의 비율은 함수에 따라 $0.05\%$에서 $2.2\%$ 사이이다. 또 커널마다 같은 IEEE 연산을 같은 순서로 하는 스칼라 짝을 두어서 원소의 값이 배열 안의 위치나 호출의 길이에 따라 달라지지 않게 했다. 곧 $f(x)$의 $i$번째 원소는 언제나 $f(x_i)$와 같은 비트이다. 넘침과 언더플로 영역, NaN, 정의역의 경계처럼 레인이 다루지 않는 인자는 C 라이브러리에 맡기며, 특수값의 규칙은 C99의 부록 F를 따른다.
\end{instardust}

\section{Boundary Cases and Special Values}
\label{sec:edge-cases}

다항식과 표만으로는 함수가 완성되지 않는다. $e^x$는 $x>\ln(\max\mathbb F)\approx709.78$이면 넘쳐 $+\infty$가 되어야 하고, $x$가 아주 작은 음수이면 서브노멀 수를 거쳐 $0$이 되어야 하며, $e^{\mathrm{NaN}}=\mathrm{NaN}$, $e^{-\infty}=0$이어야 한다. 이런 경계에서는 축소한 값 $2^k$ 자체가 넘치거나 서브노멀이 되므로 재구성 단계를 따로 다루어야 한다. 그래서 stardust의 \code{exp} 커널은 $|x|\le708$에서만 표의 방법을 쓰고 그 밖은 C 라이브러리에 맡긴다.

\code{expm1}은 $x$가 $0$ 근처일 때 또 다른 문제를 만난다. 표의 꼴 $s(1+t)-1=(s-1)+st$에서 $s-1$과 $st$가 서로 지워지기 때문이다. 그래서 $|x|<2^{-5}$에서는 $x^8$까지의 Taylor 다항식을 FMA 하나로 $x$에 더하고, $|x|<2^{-54}$이면 $e^x-1$을 반올림한 값이 $x$ 자체이므로 $x$를 그대로 낸다. 같은 함수가 인자의 범위에 따라 서로 다른 알고리즘으로 계산되는 것은 흔한 일이며, 각 범위에서 어느 방법이 가장 정확한지를 따져서 경계를 정한다.

\section{Concluding Remarks}

초월함수를 만드는 일은 이 책의 여러 장이 만나는 곳이다. 인자 축소는 8장의 비트 배치와 Sterbenz의 보조정리, 10장의 오차 없는 변환에 기대고, 다항식 근사는 Chebyshev의 이론에, 다항식 계산은 9장의 오차 해석과 13장의 컴퓨터 구조에 기댄다. 39장부터 다룰 특수함수들도 같은 틀을 따르는데, 다만 함수의 성질이 훨씬 다양하므로 인자 축소 대신 3장의 점근 전개, 6장의 점화식, 연분수, 적분 표현을 구간마다 골라 쓴다.

초월함수 구현의 이론과 실제는 Muller의 책\cite{muller2016elementary}이 표준적인 참고서이고, 근사 이론은 Trefethen의 책\cite{trefethen2019atap}이 현대적이고 친절하다.

\chapterreferences
